\documentclass[10pt,a4paper]{amsart}
\usepackage[margin=19mm]{geometry}
\usepackage{amsmath,amssymb,mathtools}
\usepackage[scr=rsfs,cal=euler]{mathalfa}
\usepackage{xfrac}
\usepackage[unicode,hidelinks,bookmarks=false]{hyperref}
\newcommand{\ie}{i.e.\ }
\newcommand{\cf}{cf.\ }
\newcommand{\resp}{respectively\ }
\newcommand{\Subsection}{\subsection}
\providecommand{\nobreakdash}{\nobreak\hbox{-}}
\setlength{\emergencystretch}{2em}
\multlinegap0pt
\usepackage{bm}
\usepackage{bbm}
\usepackage{dsfont}
\usepackage{xspace}
\usepackage{cancel}
\usepackage{mathtools}
\usepackage{amsthm}
\makeatletter
\@ifundefined{theorem}{\newtheorem{theorem}{Theorem}}{}
\@ifundefined{proposition}{\newtheorem{proposition}[theorem]{Proposition}}{}
\makeatother
\DeclareMathOperator{\rank}{rank}
\newcommand{\Gr}{G}
\newcommand{\NumWild}{406}
\newcommand{\QQ}{\mathbb{Q}}
\newcommand{\WildMinMargin}{1.1\times10^{-3}}
\newcommand{\WildRootBox}{19}
\title[Compact Hyperbolic Coxeter Six-dimensional Polytopes With Te]{Compact Hyperbolic Coxeter Six-dimensional Polytopes With Ten Facets}
\author{John Mcleod}
\date{Source: arXiv:2608.02894v1 (CC BY 4.0)}
\begin{document}
\maketitle

\noindent\emph{Source and licence.} Compact Hyperbolic Coxeter Six-dimensional Polytopes With Ten Facets. Version arXiv:2608.02894v1 (CC BY 4.0), distributed under the Creative Commons Attribution 4.0 International licence, as recorded in the arXiv licence metadata for this exact version and shown on the abstract page https://arxiv.org/abs/2608.02894v1. Full rights notice: licenses/arxiv-2608.02894-CC-BY-4.0.txt.\par\smallskip

\noindent\textit{Editorial scope.} Counted unit: the Proposition whose source label is \texttt{prop:wildexact} in the article (source0.tex lines 1116--1121, source label \texttt{prop:wildexact}), delivered together with the article's own complete original proof of it (source0.tex lines 1123--1167, one \texttt{proof} environment of 45 lines). No mathematics is written, repaired or re-derived: statement and proof are the article's own text line for line. Required context retained uncounted: none. Every definition, display and label of those lines is retained unchanged. Omitted from this package and not redistributed: 1--1115, 1122--1122, 1168--1826. Nothing inside a retained excerpt is omitted or altered: every retained body is delivered exactly as the article's lines are written, so no omission receipt of any kind accompanies this package. Citations: the retained excerpts carry no citation command at all, so no bibliography entry of the article is transcribed and no reference is redistributed. Reference adaptation: none was needed and none was made; every reference inside the delivered unit resolves inside it, so no unresolved reference or citation is produced. Printed slips kept literal: no printed word, symbol, hypothesis or number of the counted statement or of its proof was altered. Layout and class: the article's own class is \texttt{amsart} with packages including fontenc, inputenc, amsmath, amssymb, amsthm, mathtools, booktabs, longtable, tikz, array, microtype, geometry, hyperref; the wrapper is the lane's pinned \texttt{amsart} editorial wrapper at 10pt on A4, which restates the theorem-like environments the retained text needs and adds the article's own macro definitions that the retained text needs (\texttt{\textbackslash Gr}, \texttt{\textbackslash NumWild}, \texttt{\textbackslash QQ}, \texttt{\textbackslash WildMinMargin}, \texttt{\textbackslash WildRootBox}, \texttt{\textbackslash rank}), with extra wrapper packages (bm, bbm, dsfont, xspace, cancel, mathtools). The delivered numbering is the unit's own. Assets: no retained span contains a figure environment, a graphics-inclusion command, a PDF-page inclusion, a table or a TikZ picture; no image file is copied into this package, the package declares no asset dependency and no image file is redistributed with it. Licence: version arXiv:2608.02894v1 of this article is distributed under the Creative Commons Attribution 4.0 International licence, as recorded in the arXiv licence metadata for this exact version and shown on the abstract page https://arxiv.org/abs/2608.02894v1. A full rights notice is delivered at licenses/arxiv-2608.02894-CC-BY-4.0.txt. Characters: every retained excerpt is pure ASCII, so the delivered engine, XeLaTeX with its default Latin Modern text font, reports no missing character.
\begin{proposition}
\label{prop:wildexact}
Each of the $\NumWild$ wildcard-bearing labellings is infeasible, and this is
certified exactly: no local search, no bound on the ultraparallel weights, and no
floating-point tolerance enters the refutation.
\end{proposition}

\begin{proof}
Fix such a labelling, with dashed edges $D$ and wild edges $W$. Realizability
requires $\rank\Gr\le d+1$, hence the vanishing of \emph{every} $(d+2)\times(d+2)$
principal minor of $\Gr$, so exhibiting one principal minor that vanishes nowhere
on the domain refutes the labelling.

First compactify. Substituting $x_e = 1/t_e$ maps $x_e\in(1,\infty)$ to
$t_e\in(0,1)$. Each unknown occupies exactly the two symmetric positions of one
edge, so a principal minor is a polynomial of degree at most $2$ in each $x_e$ and
in each $c_e$; multiplying the minor on the index set $S$ by $t_e^{2}$ for each
$e\in D$ with both endpoints in $S$ therefore turns it into a polynomial in the
variables $(t_e)_{e\in D}$, $(c_e)_{e\in W}$. On the domain $t_e>0$, so this
changes no zero there; on the closed box it can only \emph{add} zeros, at
$t_e=0$, which makes a certificate harder to obtain and never makes one unsound.
The domain is contained in the closed rational box
\[
t_e\in[0,1]\quad(e\in D),\qquad c_e\in[c_7,1]\quad(e\in W),
\]
where $c_7$ is any rational lower bound for $\cos(\pi/7)$. Both enlargements ---
closing the box and lowering $\cos(\pi/7)$ to a rational --- only \emph{grow} the
domain, so emptiness on the box implies emptiness on the domain. The unbounded
direction has been removed rather than truncated, and no bound on $x_e$ is
assumed.

Next, the coefficients are exact. The ordinary entries $-\cos(\pi/m)$ for
$m\in\{2,\dots,6\}$ lie in $K=\QQ(\sqrt2,\sqrt3,\sqrt5)$, hence so do the
coefficients of every minor. For a minor with $r$ unknowns, which has degree at
most $2$ in each, those coefficients are recovered by tensor-product interpolation
from $3^{r}$ determinant evaluations over $K$ at integer nodes; $K$ is a field of
degree $8$ over $\QQ$ with the basis $\{\sqrt{\prod S}: S\subseteq\{2,3,5\}\}$, and
all of this is rational arithmetic on $8$-tuples.

Finally, emptiness on the box is certified by interval arithmetic. Enclosing each
coefficient by rationals and evaluating with \emph{outward}-rounded rational
intervals yields an enclosure of the polynomial's range over the box that is
guaranteed to be a superset of the true range; if that enclosure excludes $0$, the
box contains no zero, and the conclusion is a proof rather than an estimate. In
each of the $\NumWild$ cases a single $(d+2)$-principal minor settles the matter on
the root box, with no subdivision required. All but $\WildRootBox$ of them are
settled without interval arithmetic at all: the minor's index set contains neither a
dashed nor a wild edge, so the minor is a nonzero element of $K$ and the ordinary
labels alone force $\rank\Gr\ge d+2$. For the remaining $\WildRootBox$ the minor's
range over the whole box excludes $0$, in every case by a margin of at least
$\WildMinMargin$.
\end{proof}

\end{document}
